%%%%%%%%%%%%%%%%%%%%%%%%%%%%%%%%%%%%%%%%%%%%%%%%%%%%%%%%%%%%%%%%%%%%%%%%%%%%%%%
%% Phase 4 Internal Reproduction Report
%% Reproducing Huang+2021 "shielded" DESI Legacy lens finder + DR8 deployment
%% from paper + public code only.
%%
%% Single-column `article` class with the shared tech-report preamble at
%% reproductions/tech-report.sty. Builds with vanilla TeX Live (>= 2022) via
%% the Makefile in this directory.
%%%%%%%%%%%%%%%%%%%%%%%%%%%%%%%%%%%%%%%%%%%%%%%%%%%%%%%%%%%%%%%%%%%%%%%%%%%%%%%
\documentclass[11pt,letterpaper]{article}
\usepackage{../../tech-report}

\newcommand{\ResNet}{\textsc{ResNet}}
\newcommand{\Lanusse}{\textsc{CMU DeepLens}}
\newcommand{\Shield}{\textsc{shield}}
\newcommand{\Adam}{\textsc{Adam}}

\title{Reproducing the Huang 2021 Shielded ResNet\\
       and its DESI Legacy DR8 Deployment}
\author{Greg Benson \and Xiaosheng Huang}
\date{May 2026}

\begin{document}

\maketitle
\subtitle{Phase 4 Internal Reproduction Report}

\begin{abstract}
\textcite{huang2021desi} extends the \textcite{huang2020decals} DECaLS lens
search in three ways: a new ``shielded'' \ResNet{} architecture that inserts
$1{\times}1$ convolutional layers between every three residual blocks of the
\textcite{lanusse2018cmu} backbone, a deployment over the full $\sim$14{,}000~deg$^2$
DR8 footprint (adding the northern BASS/MzLS sky the DR7 search lacked), and a
two-model ensemble (original L18 + shielded) scored at a low $p\ge0.1$ threshold,
yielding $1{,}312$ published candidates. We reproduce the shielding architecture
from the methods section alone: our reconstruction has $59{,}905$ trainable
parameters (a $58.6\times$ reduction from the L18 backbone's $3{,}508{,}833$),
$4$ shields, and the full $15$-block depth, reproducing the paper's
``$\sim$60 thousand parameters'' and ``$\sim$50$\times$ reduction.'' Trained on
the identical cutouts/seed/split as our L18 baseline, the shielded net matches
the L18 held-out \AUC{} to within $\pm0.002$ ($0.9988$ vs $0.9991$ on DR9;
$0.9955$ vs $0.9943$ on DR7) --- a faithful reproduction of the paper's central
claim that shielding removes $\sim$98\% of parameters \emph{without} degrading
performance. We rebuild the published $1{,}312$-candidate catalog from the
NeuraLens release, recovering the paper's grade breakdown exactly
($216$~A~$+$~$199$~B~$+$~$897$~C). We then deploy both models on the full DR8
footprint --- $17.3$~M DEV/COMP galaxies over both the south DECaLS and north
BASS/MzLS coadds, scored in $\sim$16~h on two L4 GPUs via footprint-aware
per-brick inference. A domain-shift trap surfaced and was fixed: the
southern-trained models, having seen northern lens positives but \emph{no}
northern negatives, over-fired badly on BASS/MzLS imaging ($91\%$ of north
non-lenses scored $\ge0.1$ for L18 vs $4.8\%$ in the south); re-training with
northern negatives collapsed this to $0.8\%$, and the $60$K-parameter shielded
network proved markedly more robust to the shift than the $3.5$~M L18. Against
the $1{,}312$ published candidates, the leakage-aware analysis cleanly separates
the circular result ($86\%$ recovery at $p\ge0.9$ of the $949$ training
positives) from the honest one ($50\%$ recovery of the $363$ never-trained
shielded discoveries). \S\ref{sec:caveats} discusses the methodological
substitutions, which (as in Phase 3) bias our absolute \AUC{}s upward relative
to the paper's $0.992/0.997$.
\end{abstract}

\tableofcontents
\bigskip

% =============================================================================
\section{Introduction}
% =============================================================================

This report documents Phase~4 of an 8-phase reproduction roadmap covering
Prof.\ Xiaosheng Huang's strong-gravitational-lensing program. Phase~1
(\textcite{huang2025foundry1}; Foundry~I), Phase~2 \autocite{hsu2025pairwise},
and Phase~3 (\textcite{huang2020decals}; the DECaLS DR7 lens finder, in Phases
3a training and 3b full deployment) are reproduced in companion
tech-reports.\footnote{\texttt{reproductions/\{foundry-i,hsu-2025,huang-2020\}/papers/main.tex}.}
Phase~4 reproduces \textcite{huang2021desi}, the immediate successor to
\textcite{huang2020decals} in the lens-finder lineage.

\textcite{huang2021desi} is not a re-run of the 2020 search with more sky; it
introduces a genuinely new architecture (\S\ref{sec:shield}) and was the first
of the Huang-group searches to cover the northern BASS/MzLS footprint. As with
all prior phases, the reproduction was performed without access to Huang's
training code, training set, or trained weights. We split the work into
\textbf{Phase~4a} (the shielded architecture, validated by a controlled
comparison against our L18 baseline; \S\ref{sec:shield}--\ref{sec:compare}) and
\textbf{Phase~4b/4c} (the DR8 two-model deployment and recovery of the published
candidates; \S\ref{sec:deploy}--\ref{sec:recovery}).

% =============================================================================
\section{The shielding architecture}
\label{sec:shield}
% =============================================================================

The backbone is the \Lanusse{} \ResNet{} of \textcite{lanusse2018cmu} reproduced
in Phase~3a: a $7{\times}7$ stem followed by five stages of three pre-activated
bottleneck blocks \autocite{he2016preact}, $15$ blocks in total, with channel
progression $32\to32\to64\to128\to256\to512$ and a strided downsample at the
first block of stages~2--5. At $101{\times}101{\times}3$ input this is
$3{,}508{,}833$ parameters, of which the wide $256/512$-channel final stage
alone accounts for $\sim$75\%.

\textcite{huang2021desi} (\S3.3) modify this with ``shielding'' layers, quoting:
\begin{quote}\small
These `shields' are $1{\times}1$ convolutional layers inserted between every
three blocks of the L18 architecture\ldots they have the effect of reducing
dimensionality\ldots we reduce the number of trainable parameters by a factor of
50 (from 3 million to 60 thousand)\ldots the validation \AUC{} has increased
from 0.992\ldots to 0.997. \ldots in the final block\ldots we experimented with
reducing the output from 512 channels to 256, 128, 64, 32, and 16 channels. We
find that `shields' that keep the output to 32 channels perform the best.
\end{quote}
The rationale is regularisation by dimensionality compression: telling lenses
from non-lenses, though hard, needs only a low-dimensional latent space, so
shields ``[force] the network to learn only the most salient features''
(inspired by the InceptionNet $1{\times}1$ bottlenecks of
\textcite{szegedy2015inception}).

\paragraph{Reconstruction.} The paper does not publish the exact intermediate
channel counts, so we reconstruct (lives at
\texttt{01b\_shielded\_resnet.py}): a pre-activated shield
(\textsc{BN}~$\to$~ELU~$\to$~$1{\times}1$~conv) is inserted after each of the
\emph{first four} three-block stages --- four shields, ``between every three
blocks'' --- each squeezing the channel count to $16$, while the per-stage block
width is capped at $32$ output channels and the final stage outputs
$\mathtt{final\_out}$ channels. The L18 depth ($15$ blocks) and spatial
downsampling schedule are preserved; only the wide final stages, which dominate
the L18 parameter count, are collapsed. The swept channel ladder reproduces the
paper's experiment:

\begin{center}\small
\begin{tabular}{r r}
\toprule
\texttt{final\_out} & parameters \\
\midrule
512 & 2{,}493{,}025 \\
256 & 664{,}929 \\
128 & 205{,}537 \\
64  & 89{,}505 \\
\textbf{32} & \textbf{59{,}905} \\
16  & 58{,}257 \\
\bottomrule
\end{tabular}
\end{center}

At the paper's preferred $\mathtt{final\_out}=32$ the model has $59{,}905$
parameters --- matching the paper's ``$\sim$60 thousand'' --- a $58.6\times$
reduction from the L18 backbone. A smoke test asserts the param count, exactly
four shields, $15$ bottleneck blocks, and a $\ge50\times$ reduction
(figure~\ref{fig:shieldarch}).

\begin{figure}[ht]
\centering
\includegraphics[width=\linewidth]{figures/shielded_arch.png}
\caption{The reconstructed shielded \ResNet{}: stem, then
$[\text{stage}\times3 \to \Shield]\times4$, then a final 3-block stage and the
pooled $\mathrm{FC}(32\to1)$ head. Four $1{\times}1$ shields, $\sim$60K
parameters.}
\label{fig:shieldarch}
\end{figure}

% =============================================================================
\section{Controlled architecture comparison (Phase 4a)}
\label{sec:compare}
% =============================================================================

To isolate the architecture as the only variable we train the shielded net on
\emph{exactly} the inputs of our Phase~3a/3b L18 reproduction --- the same
$949$ NeuraLens L18-model positives, the same random DR1 negatives, the same
$\mathrm{seed}=2026$, the same $70/20/10$ stratified split --- for both the DR9
and the (paper-exact) DR7 cutout sets. This mirrors the paper's own
same-data L18-vs-shielded comparison. Training is the unchanged Lanusse recipe
(\Adam, $\mathrm{lr}=10^{-3}$ decayed $/10$ every $40$ epochs, batch $128$, $120$
epochs); the $60$K-parameter model trains in $\sim$15~min on one L4.

\begin{table}[ht]
\centering
\begin{tabular}{l r r r}
\toprule
Run & parameters & val \AUC{} & test \AUC{} \\
\midrule
L18 / DR9            & 3{,}508{,}833 & 0.9983 & 0.9991 \\
L18 / DR7            & 3{,}508{,}833 & 0.9890 & 0.9943 \\
\textbf{shielded / DR9} & \textbf{59{,}905} & \textbf{0.9989} & \textbf{0.9988} \\
\textbf{shielded / DR7} & \textbf{59{,}905} & \textbf{0.9875} & \textbf{0.9955} \\
\bottomrule
\end{tabular}
\caption{Phase~4a controlled comparison. The shielded net matches the L18
baseline to within $\pm0.002$ \AUC{} at $58.6\times$ fewer parameters. On DR9
the shielded \emph{validation} \AUC{} edges L18 ($+0.0006$); on DR7 the shielded
\emph{test} \AUC{} edges L18 ($+0.0012$). Built by
\texttt{06\_compare\_architectures.py}.}
\label{tab:arch}
\end{table}

Table~\ref{tab:arch} is the headline Phase~4a result. The differences between
the two architectures are within run-to-run noise ($\pm0.001$--$0.002$ \AUC{}),
so our reproduction confirms the paper's qualitative claim --- shielding removes
$\sim$98\% of the parameters without degrading discrimination --- even though we
do not reproduce the paper's specific $0.992\to0.997$ \emph{increase}. The
direction of the gap is consistent: in our leakage-affected setup both
architectures already sit at \AUC{}~$\approx0.99$ (see \S\ref{sec:caveats}),
leaving little headroom for shielding to demonstrate an improvement. The honest
comparison is the DR7-trained pair, where the L18 baseline is not inflated by
the DR9 NeuraLens overlap; there the shielded net is marginally \emph{better} on
the held-out test set. Figures~\ref{fig:archcmp} show the parameter-vs-\AUC{}
trade-off and the validation-\AUC{} training curves.

\begin{figure}[ht]
\centering
\includegraphics[width=0.49\linewidth]{figures/arch_comparison.png}
\hfill
\includegraphics[width=0.49\linewidth]{figures/arch_training_curves.png}
\caption{Left: held-out test \AUC{} vs.\ trainable parameters (log scale); the
shielded points sit $\sim$60$\times$ to the left of the L18 points at the same
\AUC{}. Right: validation \AUC{} per epoch, shielded (DR9, DR7) vs.\ the L18 DR7
baseline.}
\label{fig:archcmp}
\end{figure}

% =============================================================================
\section{The published catalog}
\label{sec:catalog}
% =============================================================================

\textcite{huang2021desi} publish $1{,}312$ candidates via the NeuraLens
release.\footnote{\url{https://sites.google.com/usfca.edu/neuralens}} The
release table (on disk as \texttt{neuralens\_catalog.csv}) carries, per
candidate, the \texttt{DESI-RA$\pm$DEC} name, the averaged visual grade
(Score $\in\{2.0,2.5,3.0,3.5,4.0\}$), the discovering model (\texttt{L18} or
\texttt{shielded}), and the survey region (\texttt{DECaLS} or \texttt{MzLS}).
Applying the paper's letter-grade mapping (\S3.5: A~$\ge3.5$, B~$=3.0$,
C~$\in\{2.0,2.5\}$) in \texttt{13\_extract\_huang2021\_catalog.py} reproduces
Table~3 of \textcite{huang2021desi} exactly:

\begin{table}[ht]
\centering
\begin{tabular}{l r r r r}
\toprule
Grade & L18 model & shielded model & total & paper Table 3 \\
\midrule
A   & 159 & 57  & 216 & 216 \\
B   & 139 & 60  & 199 & 199 \\
C   & 651 & 246 & 897 & 897 \\
\midrule
Total & 949 & 363 & 1{,}312 & 1{,}312 \\
\bottomrule
\end{tabular}
\caption{Published Huang+2021 candidate catalog, reconstructed from the
NeuraLens release. The grade counts, the per-model split, and the total all
match the paper exactly. $185$ of the $1{,}312$ lie in the northern MzLS region
--- the footprint Phase~3b's DR7 search could not reach.}
\label{tab:catalog}
\end{table}

This catalog is the recovery target for Phase~4b/4c (\S\ref{sec:recovery}). Note
that the $949$ L18-model rows are precisely the positives our networks were
trained on (Phase~3a), so they constitute training-set leakage; the $363$
shielded-model rows were never in training and are the leak-free test set.

% =============================================================================
\section{DR8 two-model deployment (Phase 4b)}
\label{sec:deploy}
% =============================================================================

\paragraph{Parent sample.} \textcite{huang2021desi} score every DR8 galaxy with
$\ge3$ passes in each of $grz$, $z$-band magnitude $<20.0$, typed DEV or COMP
(the REX type is a secondary pass). We reproduce these cuts over both DR8 sweep
regions --- \texttt{dr8/south/sweep/8.0/} ($437$ files, DECaLS/DECam) and
\texttt{dr8/north/sweep/8.0/} ($286$ files, BASS+MzLS) --- in
\texttt{10\_select\_parent\_sample\_dr8.py}, tagging each surviving row with a
\texttt{footprint} column (south~$=$~\texttt{RELEASE}~8000, north~$=$~8001).
This yields \textbf{17{,}290{,}814} DEV/COMP galaxies ($12.27$~M south $+$ $5.02$~M
north), $12\%$ above the paper's $\sim$15.4~M --- the same mild over-inclusivity
as our DR7 parent sample ($+9\%$), from slightly looser type/flux handling.

\paragraph{Footprint-aware brick inference.} We reuse the Phase~3b engineering
pivot --- per-brick coadd FITS downloads with local WCS slicing, $\sim$200$\times$
faster than the rate-limited cutout endpoint --- generalised in
\texttt{11b\_brick\_inference\_dr8.py} to (i) route each brick's $grz$ image
download to the correct \texttt{dr8/\{south,north\}/coadd/} directory using the
row's footprint, and (ii) score every brick with \emph{both} the L18 and
shielded checkpoints in a single download pass (network being the bottleneck),
each model normalised with its own stored band statistics. Work is partitioned
into shards by a hash of the \texttt{(footprint, brickname)} key and run across
the two L4 GPUs. The north footprint is the highest-risk new code path; a
gate (\texttt{08\_smoketest\_dr8.py}) verifies that known lenses in both the
south (DECaLS) and north (MzLS) regions download, slice, and score correctly
before the full run.\footnote{North DR8 coadds are only $\sim$50\% $grz$-complete
at the brick level, but parent-sample galaxies (which require $\mathrm{NOBS}_z\ge3$)
live only in $z$-covered, hence $grz$-complete, bricks --- so the partial bricks
never enter the deployment.}

The full sweep scored all $17{,}290{,}814$ galaxies across $298{,}843$
(footprint,~brick) units with zero brick failures, on the two L4s in
$\sim$16~h (four shards, two per GPU).

\paragraph{A domain-shift trap: north needs north negatives.} Our first
deployment used the Phase~4a checkpoints, trained on southern (DECaLS/DES)
positives and \emph{southern} DR1 negatives. They had seen $\sim$160 northern
(MzLS) lens \emph{positives} but \emph{zero} northern \emph{negatives}, leaving
the decision boundary in BASS/MzLS imaging space undefined. The consequence was
stark and footprint-localised: on the northern DR8 footprint the L18 model
\emph{over-fired} --- $91\%$ of non-lens galaxies scored $\ge0.1$ (median
$0.45$) versus $4.8\%$ in the south --- while the shielded model was far more
robust ($10\%$ vs $0.8\%$ after the fix; see below), a direct dividend of its
low-dimensional latent space. We fixed this faithfully (the route Huang trained
across the full footprint): we added $154$ northern positives (re-fetched at
DR8) and $2{,}883$ random northern negatives (\texttt{18\_build\_north\_train\_cutouts.py})
and retrained both networks (\texttt{05c\_train\_northaug.py}). The northern
over-firing collapsed --- L18 $\ge0.1$ from $91\%\to0.8\%$, shielded
$10\%\to1.4\%$ --- with the south unchanged and in-domain test \AUC{} preserved
($0.9985$ L18, $0.9996$ shielded). The deployment below uses these
north-augmented checkpoints; the merged score distributions are well-behaved
($2.9\%$/$3.8\%$ of all $17.3$~M galaxies score $\ge0.1$ for L18/shielded).
This is itself a reproduction finding: a southern-trained lens finder does
\emph{not} transfer to BASS/MzLS imaging without northern negatives, and the
$60$K-parameter shielded network degrades far more gracefully than the $3.5$~M
L18 under that domain shift.

% =============================================================================
\section{Recovery and leakage analysis (Phase 4c)}
\label{sec:recovery}
% =============================================================================

\texttt{14\_crossmatch\_recovery\_dr8.py} cross-matches both models' DR8 scores
against the $1{,}312$ published candidates (5\arcsec\ radius), tagging each as
\emph{leaked} (within 5\arcsec\ of a training positive --- the $949$ L18-model
rows) or \emph{leak-free} (the $363$ shielded-model discoveries our networks
never saw). $1{,}100$ of the $1{,}312$ fall inside the parent sample (the rest
fail the $z<20$/NOBS/type cuts or lie just outside coverage, as the $7$-missing
did for DR7). Table~\ref{tab:recovery} gives the two-model-combined recovery.

\begin{table}[ht]
\centering
\begin{tabular}{l r r r r}
\toprule
Bucket & $n$ & $p\ge0.1$ & $p\ge0.5$ & $p\ge0.9$ \\
\midrule
all ($1{,}312$)        & 1312 & 0.832 & 0.818 & 0.761 \\
\;\;leaked ($949$, training positives) & 949 & 0.927 & 0.920 & \textbf{0.860} \\
\;\;\textbf{leak-free ($363$ shielded discoveries)} & 363 & 0.584 & 0.551 & \textbf{0.504} \\
\bottomrule
\end{tabular}
\caption{Recovery of the published Huang+2021 candidates by our two-model
ensemble (combined max score). The leaked vs.\ leak-free split is the crux:
$86.0\%$ of the $949$ training positives are recovered at $p\ge0.9$ (circular),
but only $50.4\%$ of the $363$ never-trained shielded discoveries --- the latter
is the honest generalisation measure. By grade, leak-free Grade-A recovery is
$57.9\%$ at $p\ge0.9$ ($68.4\%$ at $p\ge0.1$). The shielded model edges L18 on
the leak-free set ($47.7\%$ vs $44.9\%$ all-grade at $p\ge0.9$), consistent with
its better out-of-distribution behaviour.}
\label{tab:recovery}
\end{table}

The $\sim$36~pp leaked$-$honest gap at $p\ge0.9$ ($86.0\%$ vs $50.4\%$) is a
direct measurement of the test-set leakage that the absolute \AUC{}s
(\S\ref{sec:caveats}) only hint at: recovery of lenses the model trained on is
not evidence of generalisation. The honest $\sim$50\% leak-free recovery, with
no curated hard negatives and a from-scratch training set, is a reasonable
floor.

\texttt{15\_extended\_crossmatch.py} classifies the top $2{,}000$ shielded
candidates by provenance ($10\arcsec$): $9.1\%$ are NeuraLens-L18
(training leak), $10.3\%$ Huang+2021, $3.1\%$ Huang+2020, $1.2\%$
NeuraLens-shielded, $0.9\%$ Hsu+2025; $13.9\%$ match some known catalog and
$86.2\%$ are unmatched --- closely tracking the $89\%$ unmatched fraction of the
DR7 Phase~3b top-$2{,}000$. As there (and as for the original Huang searches),
the unmatched majority is dominated by false positives, not a discovery list;
\texttt{16\_build\_inspection\_viewer.py} renders Lupton-stretched $grz$
thumbnails of the top scorers per model ($\star$ for published) for the visual-grading
pass that would distinguish the two.

\begin{figure}[ht]
\centering
\includegraphics[width=0.49\linewidth]{figures/recovery_dr8_by_grade.png}
\hfill
\includegraphics[width=0.49\linewidth]{figures/recovery_dr8_two_model.png}
\caption{Left: leak-free recovery (the $363$ shielded discoveries) by grade at
$p\ge0.1$, L18 vs shielded vs combined. Right: per-candidate L18 vs shielded
scores for the leak-free set --- the two models agree on the clear cases but
populate complementary off-diagonal regions, the basis for the paper's
two-model ensemble.}
\label{fig:recovery}
\end{figure}

% =============================================================================
\section{Reproducibility caveats}
\label{sec:caveats}
% =============================================================================

\begin{enumerate}
\item \textbf{No Huang code released.} The shielded architecture, the training
set, and the deployment glue are reconstructed from the methods section.
\item \textbf{Shield channel counts are a design choice.} The paper specifies
only ``$1{\times}1$ between every three blocks,'' ``final block $32$ best,'' and
``$\sim$60K parameters.'' Our $\mathtt{SHIELD\_CH}=16$ reconstruction lands at
$59{,}905$ parameters --- the closest principled fit --- but the exact
intermediate widths may differ from Huang's.
\item \textbf{Training-set leakage.} As in Phase~3a, our positives are the
NeuraLens L18-model rows, which are themselves Huang+2021 candidates. This
inflates our absolute \AUC{}s ($\sim$0.99 vs.\ the paper's $0.992/0.997$) and
means the controlled comparison (Table~\ref{tab:arch}) is the trustworthy
result, not the absolute numbers. The leak-free recovery of the $363$
shielded-model candidates (\S\ref{sec:recovery}) is designed to side-step this.
\item \textbf{Random vs.\ curated negatives.} We use uniformly random DR1
galaxies as non-lenses, not the paper's by-eye hard negatives (spirals, rings,
mergers, artefacts), which were not released.
\item \textbf{GPU device ordering.} On this host (8$\times$A16~$+$~2$\times$L4),
CUDA's default \texttt{FASTEST\_FIRST} ordering does not match \texttt{nvidia-smi}
indices; \texttt{CUDA\_DEVICE\_ORDER=PCI\_BUS\_ID} is required to target the two
L4s.
\end{enumerate}

% =============================================================================
\section{Software and data availability}
% =============================================================================

All code lives at \texttt{reproductions/huang-2021/} (scripts \texttt{01b},
\texttt{05}--\texttt{16}); the L18 backbone (\texttt{01\_lanusse\_resnet.py}) and
all training inputs are shared with \texttt{huang-2020/} via symlink. The
shielded architecture is \texttt{01b\_shielded\_resnet.py}; the controlled
comparison is reproduced by \texttt{06\_compare\_architectures.py}. DR8 sweep
catalogs ($\sim$1.5~TB) and per-brick coadds are streamed from the NERSC Legacy
Surveys portal and not redistributed. The reconstructed published catalog
(\texttt{data/huang2021\_published\_catalog.csv}) and the architecture-comparison
table (\texttt{data/arch\_comparison.csv}) are committed.

% =============================================================================
\section{Conclusions}
% =============================================================================

Phase~4a reproduces the central methodological contribution of
\textcite{huang2021desi}: a shielded \ResNet{} with $58.6\times$ fewer
parameters ($59{,}905$ vs $3{,}508{,}833$) that matches the L18 backbone's
held-out \AUC{} to within $\pm0.002$ on identical data. We reconstruct the
published $1{,}312$-candidate catalog exactly, including its grade and
per-model breakdown and its $185$ northern MzLS candidates. Phase~4b/4c then
deploys both models on the full $17.3$~M-galaxy DR8 footprint --- the first
Huang-group search this reproduction extends into the northern BASS/MzLS sky.
Two results stand out: (i) a southern-trained finder needs northern negatives
to avoid badly over-firing on BASS/MzLS imaging ($91\%\to0.8\%$ after the fix),
and the shielded network degrades far more gracefully under that domain shift
than the L18 --- a second, independent vindication of the shielding idea beyond
the parameter count; and (ii) the leakage-aware recovery cleanly separates the
circular $86\%$ recovery of training positives from the honest $\sim$50\%
recovery of the $363$ never-trained shielded discoveries. The natural Phase~5
is the DR9/DR10 successors (Storfer~2024, Inchausti~2025), which add a second
backbone (EfficientNetV2) and a meta-learner on top of the shielded \ResNet{}
reproduced here.

\bibliographystyle{plainnat}
\bibliography{references}

\end{document}
